% Default Physics Experiment Report Template
% AutoReport - Built-in template
%
% Compilation: xelatex -interaction=nonstopmode main.tex
%              xelatex -interaction=nonstopmode main.tex
% (Must compile twice for cross-references)

\documentclass[12pt,a4paper]{article}

%=== Packages ===%
\usepackage{ctex}           % Chinese support
\usepackage{amsmath,amssymb} % Math symbols
\usepackage{graphicx}       % Images
\usepackage{hyperref}       % Links
\usepackage{geometry}       % Page margins
\usepackage{booktabs}       % Professional tables
\usepackage{physics}        % Physics notation

%=== Page Layout ===%
\geometry{
    left=2.5cm,
    right=2.5cm,
    top=2.5cm,
    bottom=2.5cm
}

%=== Hyperref Setup ===%
\hypersetup{
    colorlinks=true,
    linkcolor=blue,
    filecolor=magenta,
    urlcolor=cyan,
    citecolor=blue,
}

%=== Document Metadata ===%
\title{Physics Experiment Report}
\author{Student Name}
\date{\today}

\begin{document}

%=== Title Page ===%
\maketitle

%=== Sections ===%

\section{Introduction}
\label{sec:introduction}

This experiment investigates [brief description of the physical phenomenon]. The study of [topic] is fundamental to understanding [broader context]. This report presents the theoretical framework, experimental methods, data analysis, and conclusions.

The objectives of this experiment are:
\begin{enumerate}
    \item [First objective]
    \item [Second objective]
    \item [Third objective]
\end{enumerate}

\section{Theory}
\label{sec:theory}

[Begin with explanatory text before equations]

The fundamental principle governing this experiment is [principle name]. For [system description], the relevant equation is:

$$
\text{[Equation]}
$$

where [variable definitions].

[Derivation steps and explanations]

\section{Experimental Setup}
\label{sec:setup}

[Begin with overview of the apparatus]

The experimental apparatus consists of [description]. Figure~\ref{fig:setup} shows the experimental configuration.

\begin{figure}[h]
\centering
% \includegraphics[width=0.8\textwidth]{figures/setup.png}
\caption{Experimental setup. [Description of what the figure shows]}
\label{fig:setup}
\end{figure}

The procedure for data collection was as follows:
\begin{enumerate}
    \item [Step 1]
    \item [Step 2]
    \item [Step 3]
\end{enumerate}

\section{Data Analysis}
\label{sec:analysis}

[Begin with explanation of analysis approach]

The data were analyzed using [method]. For each measurement, [description of calculations].

[Equations used in analysis]

$$
\text{[Analysis equation]}
$$

The uncertainties were propagated using [method].

\section{Results}
\label{sec:results}

[Begin with overview of findings]

Table~\ref{tab:results} presents the measurements obtained from the experiment.

\begin{table}[h]
\centering
\caption{Measurement results}
\label{tab:results}
\begin{tabular}{ccc}
\toprule
Trial & Measured Value (unit) & Uncertainty (unit) \\
\midrule
1 & [value] & [uncertainty] \\
2 & [value] & [uncertainty] \\
3 & [value] & [uncertainty] \\
\bottomrule
\end{tabular}
\end{table}

Figure~\ref{fig:results} shows [description of what the figure displays].

\begin{figure}[h]
\centering
% \includegraphics[width=0.8\textwidth]{figures/results.png}
\caption{[Caption describing what the figure shows and interpreting the results]}
\label{fig:results}
\end{figure}

The measured value of [quantity] is:

$$
\text{[Result]} \pm \text{[Uncertainty]} \text{ units}
$$

\section{Discussion}
\label{sec:discussion}

[Begin with interpretation of results]

The results show that [interpretation]. Comparing with the theoretical prediction from Section~\ref{sec:theory}, the measured value agrees/disagrees within experimental uncertainty.

The observed [phenomenon] can be explained by [explanation]. Possible sources of error include:
\begin{itemize}
    \item [Error source 1]
    \item [Error source 2]
    \item [Error source 3]
\end{itemize}

\section{Conclusion}
\label{sec:conclusion}

[Begin with summary of main findings]

This experiment investigated [topic]. The key findings are:
\begin{enumerate}
    \item [Finding 1]
    \item [Finding 2]
    \item [Finding 3]
\end{enumerate}

The results confirm/support [conclusion]. Suggestions for improving the experiment include [recommendations].

\section{References}
\label{sec:references}

\begin{thebibliography}{9}

\bibitem{ref1}
[Author]. (\textit{Year}). \textit{[Title]}. [Journal], [Volume](Issue), [Pages].

\bibitem{ref2}
[Author]. (\textit{Year}). \textit{[Title]}. [Publisher].

\end{thebibliography}

\end{document}
